\centerline{\bf \Huge Теория игр}

\begin{flushright}
        \scriptsize{Когда игра заканчивается, король и пешка падают в одну и ту же коробку.\\ Итальянская пословица}
\end{flushright}

\problem{1}
На доске написаны 10 единиц и 10 двоек. За ход разрешается стереть две любые цифры и, если они были одинаковыми, написать двойку, а если разными - единицу. Если последняя оставшаяся на доске цифра - единица, то выиграл первый игрок, если двойка - то второй.

\problem{2}
Двое по очереди ставят королей в клетки доски $9 \times 9$ так, чтобы короли не били друг друга. Проигрывает тот, кто не может сделать ход.

\problem{3}
Король стоит на поле а1. За один ход его можно передвинуть на одно поле вправо, или на одно поле вверх, или на одно поле по диагонали "вправо-вверх". Выигрывает тот, кто поставит короля на поле $h 8$.

\problem{4}
Двое по очереди ставят крестики и нолики в клетки доски $9 \times 9$. Начинающий ставит крестики, его соперник - нолики. В конце подсчитывается, сколько имеется строчек и столбцов, в которых крестиков больше, чем ноликов - это очки, набранные первым игроком. Количество строчек и столбцов, где ноликов больше - очки второго. Тот из игроков, кто наберет больше очков, побеждает.

\problem{5}
Имеется две кучки камней: в первой - 7 камней, во второй - 5. За ход разрешается брать любое количество камней из одной кучки или поровну камней из обеих кучек. Проигрывает тот, кто не может сделать ход.